\pnum{5.35}{Nonrelativistic white-dwarf degeneracy pressure}
Let the star contain $N$ nucleons of mass
$M$, and $N_e=dN$ electrons of mass $m$.
Assume a uniformly dense sphere with
$V=4\pi R^3/3$. In a nonrelativistic
spin-$1/2$ Fermi gas,
\begin{align}
E_e&=\frac{3\sh^2}{10m}
(3\pi^2)^{2/3}(dN)^{5/3}V^{-2/3}\\
&=\frac{3}{10}
\left(\frac{9\pi}{4}\right)^{2/3}
\frac{\sh^2}{m}
\frac{d^{5/3}N^{5/3}}{R^2}.
\end{align}
A uniform Newtonian sphere of mass
$NM$ has gravitational potential energy
\begin{equation}
\boxed{E_g=-\frac35\frac{G(NM)^2}{R}}.
\end{equation}
The total energy has the form
$E_{\rm tot}(R)=A/R^2-B/R$, with
$A>0$, $B>0$. Differentiate:
$\dd E/\dd R=-2A/R^3+B/R^2=0$.
The equilibrium radius is $R=2A/B$:
\begin{equation}
\boxed{R=\left(\frac{9\pi}{4}\right)^{2/3}
\frac{\sh^2 d^{5/3}}{GmM^2N^{1/3}}}.
\end{equation}
At this point $\dd^2E/\dd R^2=2A/R^4>0$,
so it is an energy minimum.
Putting $d=1/2$ and conventional nucleon
masses gives approximately
$\boxed{R=7.6\times10^{25}N^{-1/3}\ {\rm m}}$.
For a solar mass $M_\odot\simeq
1.989\times10^{30}$ kg,
$N\simeq M_\odot/M\simeq1.19\times10^{57}$,
and
\begin{equation}
\boxed{R_{\rm WD}\simeq7.2\times10^3\ {\rm km}}.
\end{equation}
Compute the electron density
$n_e=dN/(4\pi R^3/3)$ and
$E_F=\sh^2(3\pi^2n_e)^{2/3}/(2m)$.
For this model the result is
$\boxed{E_F\sim2.0\times10^5\ {\rm eV}}$,
not dramatically below the electron
rest energy $m_ec^2=5.11\times10^5$ eV.
Thus relativistic corrections are significant
and the nonrelativistic white-dwarf model
is only an approximation.

\pnum{5.36}{Ultrarelativistic gas, critical stellar mass and neutron star}
(a) For $E(k)\simeq\sh ck$,
two spin orientations supply
$V k^2\dd k/\pi^2$ states.
Integrating gives
\begin{equation}
\boxed{E_{\rm ultrarel}
=\frac{V\sh c}{\pi^2}\int_0^{k_F}k^3\dd k
=\frac{V\sh ck_F^4}{4\pi^2}
=\frac34N_e\sh ck_F}.
\end{equation}
(b) With $N_e=dN$, $V=4\pi R^3/3$,
\begin{equation}
E_e^{\rm UR}=
\frac34\left(\frac{9\pi}{4}\right)^{1/3}
\frac{\sh c\,d^{4/3}N^{4/3}}R.
\end{equation}
Add $E_g=-3GM^2N^2/(5R)$. Both terms
scale as $R^{-1}$, so the
nonrelativistic $R^{-2}$ stabilization
disappears. The coefficient of $1/R$
changes sign at
\begin{equation}
\boxed{N_c=\left[
\frac54\left(\frac{9\pi}{4}\right)^{1/3}
\frac{\sh c\,d^{4/3}}{GM^2}
\right]^{3/2}}.
\end{equation}
For $d=1/2$, $M$ equal to a nucleon mass,
the estimate is
$\boxed{N_c\simeq2.04\times10^{57}}$,
equivalent to about
$\boxed{1.7M_\odot}$.
This is the uniform-density Newtonian toy
model's Chandrasekhar-like limit, not a
precision modern white-dwarf mass prediction.
(c) Suppose a solar-mass object is supported
by \emph{neutron} degeneracy rather than
electron degeneracy. Set fermion mass
$m=m_n$, gravitational nucleon mass
$M\simeq m_n$, and one fermion per nucleon
$d=1$ in the nonrelativistic radius formula:
\begin{equation}
\boxed{R_{\rm NS}\sim1.2\times10^4\ {\rm m}
\sim12\ {\rm km}}.
\end{equation}
The density and Fermi energy are
\begin{equation}
n_n=\frac{N}{4\pi R^3/3},\quad
E_{F,n}=\frac{\sh^2}{2m_n}(3\pi^2n_n)^{2/3}
\sim\boxed{6\times10^7\ {\rm eV}}.
\end{equation}
Compare with $m_nc^2\simeq939$ MeV:
the kinetic Fermi scale is roughly
$6\%$ of the rest energy, so an
initial nonrelativistic estimate is not
grossly inconsistent. In real neutron stars,
strong interactions and general relativity
are essential.

\pnum{5.37}{Density of states and one-dimensional van Hove peaks}
A periodic chain has $N$ cells of spacing
$a$, and allowed $q$ spacing
$\Delta q=2\pi/(Na)$. Counting two
momenta $\pm q$ with the same energy,
the DOS per unit length is
\begin{equation}
\boxed{g(E)\equiv\frac{G(E)}{Na}
=\frac1{\pi|\dd E/\dd q|}}.
\end{equation}
(a) For a free particle,
$E=\sh^2q^2/(2m)$, so
\begin{equation}
\boxed{g_{\rm free}(E)
=\frac1{\pi\sh}\sqrt{\frac{m}{2E}}}.
\end{equation}
(b) For the delta lattice let
$k=\sqrt{2mE}/\sh$ and
\begin{equation}
f(k)=\cos(ka)+
\frac{m\alpha}{\sh^2k}\sin(ka),
\qquad \cos(qa)=f(k).
\end{equation}
Differentiate the implicit equation:
$-a\sin(qa)\,\dd q/\dd E
=(\dd f/\dd k)(\dd k/\dd E)$,
where $\dd k/\dd E=m/(\sh^2k)$.
Since $\sin^2(qa)=1-f(k)^2$,
\begin{equation}
\boxed{g_\alpha(E)=
\frac{m}{\pi a\sh^2 k}
\frac{|f'(k)|}{\sqrt{1-f(k)^2}}},
\end{equation}
valid in allowed bands $|f(k)|<1$, with
\begin{equation}
f'(k)=-a\sin(ka)+\frac{m\alpha}{\sh^2}
\left[\frac{a\cos(ka)}k-
\frac{\sin(ka)}{k^2}\right].
\end{equation}
For $|f(k)|>1$ there are no real
Bloch modes and $g_\alpha(E)=0$.
(c) The numerical script compares
$g_0(E)$ and $g_{\alpha=1}(E)$
for $m=\sh=a=1$. Near a band edge,
the denominator $\sqrt{1-f^2}$
vanishes while the derivative is
generally nonzero, producing the
characteristic one-dimensional
van Hove divergence.
\checkit{The DOS formula counts the
two spatial momenta $\pm q$.
If electron spin is also counted,
multiply this spatial DOS by two.}

\pnum{5.38}{Harmonic chain, normal modes and phonon operators}
Consider $N$ equal masses on a ring,
$x_{N+1}=x_1$, with Hamiltonian
\begin{equation}
H=-\frac{\sh^2}{2m}\sum_{j=1}^N
\frac{\partial^2}{\partial x_j^2}
+\frac{m\omega^2}{2}
\sum_{j=1}^N(x_{j+1}-x_j)^2.
\end{equation}
(a) The geometric series satisfies
$\sum_{j=1}^N\ee^{2\pi\ii j n/N}=N$
when $n$ is divisible by $N$ and
zero otherwise. Therefore for
$1\le k,k'\le N-1$,
\begin{equation}
\frac1N\sum_j\ee^{2\pi\ii j(k-k')/N}
=\delta_{kk'},\qquad
\frac1N\sum_j\ee^{2\pi\ii j(k+k')/N}
=\delta_{k',N-k}.
\end{equation}
(b) Define $\omega_k=2\omega\sin(\pi k/N)$,
and the operators
\begin{equation}
a_k^\pm=\frac1{\sqrt N}\sum_{j=1}^N
\ee^{\pm2\pi\ii jk/N}
\left[
\sqrt{\frac{m\omega_k}{2\sh}}\,x_j
\mp\sqrt{\frac{\sh}{2m\omega_k}}\,
\partial_{x_j}\right].
\end{equation}
Using $[\partial_{x_j},x_l]=\delta_{jl}$
and the Fourier orthogonality from (a),
the only surviving terms in
$[a_k^-,a_{k'}^+]$ are the two
coordinate-derivative cross terms.
Their coefficients add to one:
\begin{equation}
\boxed{[a_k^-,a_{k'}^+]=\delta_{kk'},\qquad
[a_k^-,a_{k'}^-]=[a_k^+,a_{k'}^+]=0}.
\end{equation}
(c) Let $R=N^{-1}\sum_jx_j$.
Inverse discrete Fourier transformation,
using (a), gives
\begin{align}
\boxed{x_j}&=R+\frac1{\sqrt N}\sum_{k=1}^{N-1}
\sqrt{\frac{\sh}{2m\omega_k}}
(a_k^-+a_{N-k}^+)\ee^{2\pi\ii jk/N},\\
\boxed{\partial_{x_j}}&=
\frac1N\partial_R+
\frac1{\sqrt N}\sum_{k=1}^{N-1}
\sqrt{\frac{m\omega_k}{2\sh}}
(a_k^--a_{N-k}^+)\ee^{2\pi\ii jk/N}.
\end{align}
(d) Substituting these inverses into
the kinetic and nearest-neighbor
quadratic potential, Fourier orthogonality
eliminates all $k\ne k'$ cross terms.
The result is a free center of mass
plus $N-1$ independent oscillators:
\begin{equation}
\boxed{H=-\frac{\sh^2}{2Nm}\partial_R^2+
\sum_{k=1}^{N-1}\sh\omega_k
\left(a_k^+a_k^-+\frac12\right)}.
\end{equation}
For a center-of-mass plane wave
$\ee^{\ii K R}$ the \emph{energy is positive}:
\begin{equation}
\boxed{E=\frac{\sh^2K^2}{2Nm}+
\sum_{k=1}^{N-1}\sh\omega_k
\left(n_k+\frac12\right)},\quad
n_k=0,1,\dots.
\end{equation}
Each $n_k$ counts phonons in mode $k$.
\checkit{The sign of the free center-of-mass
energy follows from acting
$-\partial_R^2$ on $\ee^{\ii K R}$.
Any negative sign in the printed
energy eigenvalue is a typographic
inconsistency, not a negative kinetic energy.}

\pnum{5.39}{Why periodic and hard-wall free-electron gases agree}
For periodic box lengths $L_x,L_y,L_z$,
the condition
$\psi(\mathbf r+L_i\hat e_i)=\psi(\mathbf r)$
on a plane wave
$\psi=V^{-1/2}\ee^{\ii\mathbf k\cdot\mathbf r}$
forces $\ee^{\ii k_iL_i}=1$.
Therefore
\begin{equation}
\boxed{k_iL_i=2\pi n_i,\qquad
n_i\in\mathbb Z}.
\end{equation}
Each allowed $\mathbf k$ occupies
$\boxed{\Delta k_x\Delta k_y\Delta k_z
=(2\pi)^3/V=8\pi^3/V}$.
This cell is eight times larger
than the $\pi^3/V$ cell of hard-wall
standing-wave states, but periodic
$\mathbf k$ spans the \emph{entire}
sphere rather than only its positive
octant. These factors of eight cancel.
Including the same factor two for spin,
\begin{equation}
N_e=\frac{2V}{(2\pi)^3}
\frac{4\pi k_F^3}{3}
=\frac{Vk_F^3}{3\pi^2},
\end{equation}
so
\begin{equation}
\boxed{k_F=(3\pi^2n)^{1/3},\qquad
E_F=\frac{\sh^2}{2m}(3\pi^2n)^{2/3},\qquad
E_{\rm tot}=\frac35N_eE_F}.
\end{equation}
Bulk thermodynamic quantities therefore
agree in the large-$N$ limit; surface
and finite-size spectra differ.
